\section{Separation and compactness}\label{app:foundations}

Here $X$ is a general real normed space, as in Lemma~\ref{lem:weak-facts}.

\begin{proof}[Proof of Lemma~\ref{lem:weak-facts}]
We begin with the Hahn--Banach extension argument underlying the first two assertions.
Suppose that a linear functional $\ell$ on a subspace $M$ is dominated by a
continuous sublinear function $p$ on $X$. To extend it to $M+\R b$, where
$b\notin M$, choose its value at $b$ in the interval
\[
 \sup_{a\in M}\bigl(\ell(a)-p(a-b)\bigr)
 \ \leq\ \ell(b)\ \leq\
 \inf_{c\in M}\bigl(p(c+b)-\ell(c)\bigr).
\]
This interval is nonempty because
$\ell(a)+\ell(c)\leq p(a+c)\leq p(a-b)+p(c+b)$.
Its endpoints are finite, as is seen by setting $a=0$ or $c=0$ in this
inequality. The displayed bounds give domination on $a+b$ and $a-b$;
positive homogeneity then gives domination on $a+tb$ for every real $t$.
Unions of chains of dominated extensions remain dominated extensions, so
Zorn's Lemma gives an extension to $X$. It is continuous because domination
of both $v$ and $-v$ bounds its absolute value near zero.
Applying this construction with $p=\norm{\cdot}$ and
$\ell(tv)=t\norm{v}$ gives a norming functional (use the zero functional
when $v=0$). In particular,
\[
             \norm{v}=\sup_{f\in X^*,\ \norm{f}\leq1}f(v).
\]

For separation, let $q$ lie outside a nonempty norm-closed convex set $D$.
Choose $\rho>0$ such that the open convex set
$U=D+\{v:\norm{v}<\rho\}$ still excludes $q$, and fix $u\in U$.
The gauge
\[
 p(v)=\inf\{t>0:v\in t(U-u)\}
\]
is finite, nonnegative, positively homogeneous, and subadditive.
For the last property, if $v\in s(U-u)$ and $w\in t(U-u)$, convexity gives
$v+w\in(s+t)(U-u)$, and one takes infima.
Since $U-u$ is open, convex, and contains zero, $\{v:p(v)<1\}=U-u$:
one inclusion follows from convexity, and the other by extending a little
past a given point along its ray using openness.
A ball about zero contained in $U-u$ gives $p(v)\leq c\norm{v}$ for some
$c$; subadditivity then makes $p$ continuous.
On the line through $q-u$, the functional taking the value $p(q-u)$ at
$q-u$ is dominated by $p$; for negative multiples this follows from
$-p(q-u)\leq p(u-q)$. Its extension $f$ therefore satisfies
\[
             f(z-u)<1\leq f(q-u)\qquad(z\in U).
\]
In particular $f\ne0$, and taking the supremum over the ball of radius
$\rho$ added to any $d\in D$ gives
$f(d)+\rho\norm{f}\leq f(q)$. Thus $f$ separates $q$ from $D$ with a
positive gap. This proves weak closedness of $D$. Applied to closed balls
it proves weak lower semicontinuity of the norm. Separation also shows
that the weak closure of a convex set equals its norm closure.
On $X^r$, equipped with the maximum norm, every continuous linear
functional has the form $\sum_{i=1}^r f_i(x_i)$, with $f_i\in X^*$,
by restriction to the coordinate subspaces. This proves the assertion
about the finite-product topology and allows the same separation argument
there.

Here are the compactness details. A proper filter is a family of subsets
closed under finite intersections and supersets, containing the whole
space but not the empty set. By Zorn's Lemma it extends to a maximal proper
filter, called an ultrafilter. For every subset an ultrafilter contains
either that subset or its complement: if the subset cannot be added while
preserving the filter property, an existing filter member misses it.
An ultrafilter converges to a point when it contains every neighborhood
of that point. In a compact space it converges: the closures of its members
have the finite intersection property, and any point in their intersection
is a limit by the preceding either-or property. Limits are unique in a
Hausdorff space. Conversely, if every ultrafilter converges, extend any
family of closed sets with the finite intersection property to an
ultrafilter; its limit belongs to all those sets. This is the closed-set
criterion for compactness. An ultrafilter on a product of compact spaces
has a limit in each coordinate, and the tuple of these limits is its
limit, since a basic product neighborhood restricts only finitely many
coordinates. This proves product compactness.

Now let $D\subseteq X$ be bounded, norm closed, and convex, with $X$
reflexive. Given an ultrafilter $\mathcal V$ on $D$, each bounded real
evaluation has a limit in a compact interval, and
\[
             L(f)=\lim_{\mathcal V} f(d)\qquad(f\in X^*)
\]
defines a linear functional with $\abs{L(f)}\leq M\norm{f}$, where
$M=\sup_{d\in D}\norm{d}$. Linearity follows by intersecting the filter
members on which each of two evaluations is close to its limit.
Reflexivity supplies $a\in X$ with $L(f)=f(a)$ for every $f$.
Thus $\mathcal V$ converges weakly to $a$, and weak closedness gives
$a\in D$. The ultrafilter criterion proves weak compactness.
\end{proof}
